\chapter{A Tutorial}\label{ch:tutorial}

This chapter is a hands-on tour of Dylan. It starts with arithmetic and
works up, a step at a time, through variables, functions, control flow,
collections, classes, generic functions and conditions, and it ends by
building a small complete program and compiling it into an executable.
Each section introduces a few ideas with short examples that you can type
and run, and each points to the chapter in Part~II that covers the same
ground in full.

The tutorial assumes you know some other programming language, but not
Lisp or Scheme. If you do know a Lisp, you will find much that is
familiar, and you can move quickly through the first sections.

\section{How to follow along}\label{sec:tut-follow}

The easiest way to follow the tutorial is at the interactive
implementation's read-eval-print loop. Start SBCL in the top directory of
the clone and start the loop, as Chapter~\ref{ch:start} explains:

\begin{repl}
$ sbcl
* (load "src/load-dylan.lisp")
* (dylan:rep)

Entering Dylan read-eval-print-loop.
Exit by typing "dylan:done"

?
\end{repl}

Then type the examples as you come to them. The examples are written as
Dylan source, with the value of each expression shown in a comment
beginning \code{; =>}. At the loop, the same value appears after
\code{;Value:}:

\begin{dylan}
(+ 1 2)   ; => 3
\end{dylan}

When an example prints something, the text it prints is shown in a
comment beginning \code{; prints}. Every example in this chapter has been
run in both implementations: at the loop, and compiled with \cmd{dylanc}.
You can also put the examples of a section into a file and compile it, as
Section~\ref{sec:tut-compile} shows; in a compiled program, nothing is
printed unless the program prints it, so wrap an expression in
\code{print} to see its value.

\begin{notebox}
The examples build on one another. A function defined in one section is
often used in the next, so type them in order, in one session.
\end{notebox}

\section{Expressions and prefix syntax}\label{sec:tut-syntax}

\index{prefix syntax}
A Dylan program is made of \term{expressions}, also called
\term{forms}. The simplest are literal constants, which evaluate to
themselves:

\begin{dylan}
42          ; => 42
-7          ; => -7
3.5         ; => 3.5
"a string"  ; => "a string"
#\a         ; => #\a
#t          ; => #t
#f          ; => #f
\end{dylan}

\code{\#t} and \code{\#f} are the true and false values, and
\code{\#\textbackslash{}a} is the character \code{a}. Everything else is
written as a parenthesized list. A list whose first element names a
function is a \term{call}: the function comes first, and its arguments
follow, separated by white space.

\begin{dylan}
(+ 1 2)                ; => 3
(max 3 9 4)            ; => 9
(* (+ 1 2) (- 10 4))   ; => 18
\end{dylan}

This is \term{prefix syntax}: the operator comes before its operands, and
the parentheses mark exactly where each expression begins and ends. There
is no operator precedence to remember, because nesting says everything.
The last example multiplies the value of \code{(+ 1 2)} by the value of
\code{(- 10 4)}. Operators such as \code{+} and \code{max} accept any
number of arguments:

\begin{dylan}
(+ 1 2 3 4)   ; => 10
(< 1 2 3)     ; => #t
\end{dylan}

Arithmetic operators are ordinary functions, not special syntax, which is
why they are written the same way as every other call. A few lists are
not calls but \term{special forms} with rules of their own, such as
\code{if}, \code{define} and \code{method}; you will meet them in the
sections that follow.

\paragraph{Names.} Names in Dylan may contain letters, digits and many
punctuation characters, so \code{even?}, \code{set!}, \code{<point>},
\code{*count*} and \code{remove-key!} are all single names. Dylan uses
punctuation by convention to suggest what a name is for:

\begin{itemize}
\item a name in angle brackets, such as \code{<integer>}, names a class;
\item a name ending in \code{?}, such as \code{even?}, names a predicate,
  a function that answers true or false;
\item a name ending in \code{!}, such as \code{add!}, names a function
  that may modify one of its arguments;
\item a name between asterisks, such as \code{*count*}, names a module
  variable whose value changes.
\end{itemize}

Names are not case-sensitive: \code{Max}, \code{MAX} and \code{max} are the
same name. Because spaces separate the parts of a form, you must put
spaces around operators: \code{(- x 1)} subtracts, but \code{x-1} is a
single name.

\paragraph{Symbols, keywords and quoting.} A name preceded by a quote
mark is a \term{symbol}, a value in its own right rather than a reference
to a variable. A name ending in a colon is a \term{keyword}, which is used
to label arguments. Both evaluate to themselves:

\begin{dylan}
'apple       ; => apple
size:        ; => size:
\end{dylan}

The quote mark also stops a list from being evaluated, which is one way to
write a constant list. A hash sign before a list makes a constant vector:

\begin{dylan}
'(1 2 3)     ; => (1 2 3)
#(1 2 3)     ; => #(1 2 3)
\end{dylan}

\paragraph{Comments.} A semicolon starts a comment that runs to the end of
the line. Text between \code{\#|} and \code{|\#} is a comment that can span
several lines.

\smallskip
Chapter~\ref{ch:syntax} describes the syntax completely (\dylanbook{chapter
3}).

\section{Numbers and arithmetic}\label{sec:tut-numbers}

\index{numbers}
Dylan has integers, exact ratios and floating-point numbers. Division of
integers that do not divide evenly gives an exact ratio, not a rounded
result:

\begin{dylan}
(/ 10 5)      ; => 2
(/ 10 4)      ; => 5/2
(+ 1/3 1/6)   ; => 1/2
(/ 10.0 4)    ; => 2.5
(as <float> 1/4)  ; => 0.25
\end{dylan}

Mixing a float with an exact number gives a float, and \code{as} converts a
number to another class. Integer division that yields an integer is done
by the functions whose names end in a slash. Each returns two values, the
quotient and the remainder:

\begin{dylan}
(truncate/ 17 5)   ; => 3, 2
(floor/ -7 2)      ; => -4, 1
(modulo -7 2)      ; => 1
(remainder -7 2)   ; => -1
\end{dylan}

\code{truncate/} rounds the quotient toward zero and \code{floor/} rounds
it down; \code{modulo} and \code{remainder} return just the remainder that
\code{floor/} and \code{truncate/} would. Functions that return several
values are common in Dylan, and Section~\ref{sec:tut-variables} shows how
to receive them.

There are the usual comparisons and predicates:

\begin{dylan}
(= 2 2.0)     ; => #t
(< 1 2 3)     ; => #t
(even? 4)     ; => #t
(zero? 0)     ; => #t
(abs -3)      ; => 3
(min 3 9 4)   ; => 3
\end{dylan}

\code{=} compares numbers by value, whatever their class; Dylan also has
\code{id?}, which answers whether two values are the very same object.
Powers, roots and the rest of the mathematical library are there too:

\begin{dylan}
(expt 2 10)   ; => 1024
(sqrt 16)     ; => 4.0
(sqrt 2)      ; => 1.4142135623730951
\end{dylan}

Notice that \code{(sqrt 16)} answers the float \code{4.0}, not the integer
\code{4}. The 1992 manual defines \code{sqrt}, \code{exp} and the
trigonometric functions as in Common Lisp, and its own examples show them
answering floats, so both implementations do the same.

\begin{warnbox}
The two implementations differ in the size of integers. The interactive
implementation has integers of any size, as Common Lisp does. A program
compiled with \cmd{dylanc} has 63-bit integers, and an operation whose
result does not fit signals an error:
\begin{dylan}
(expt 2 62)   ; => 4611686018427387904 in the interactive implementation
              ; error: integer overflow (integers are 63 bits) when compiled
\end{dylan}
Neither implementation has complex numbers.
\end{warnbox}

Chapter~\ref{ch:numbers} covers numbers in full (\dylanbook{chapter 14}).

\section{Variables}\label{sec:tut-variables}

\index{variables}
\code{define} creates a \term{module variable}, one that the whole
program can see, and gives it a value. \code{set!} changes the value of a
variable that already exists:

\begin{dylan}
(define *count* 0)            ; => *count*
(set! *count* (+ *count* 1))  ; => 1
*count*                       ; => 1
\end{dylan}

At the loop, a \code{define} answers the name it defined, and
\code{set!} answers the new value. Writing a variable's name by itself
evaluates to its value.

\index{bind}
\code{bind} creates \term{local variables} that exist only inside its
body. Its first argument is a list of bindings, each a name and an
initial value; the remaining forms are the body, and the value of the
last one is the value of the whole \code{bind}:

\begin{dylan}
(bind ((x 3)
       (y 4))
  (sqrt (+ (* x x) (* y y))))   ; => 5.0
\end{dylan}

The bindings are made in order, so a later initial value can use an
earlier variable:

\begin{dylan}
(bind ((a 1)
       (b (+ a 1)))
  (list a b))                   ; => (1 2)
\end{dylan}

A binding with more than one name receives multiple values, one in each
name. This is how you use the quotient and remainder from
\code{truncate/}:

\begin{dylan}
(bind ((quotient rem (truncate/ 17 5)))
  (list quotient rem))          ; => (3 2)
\end{dylan}

Local variables are the usual way to name intermediate results. Module
variables are best kept for things that really belong to the whole
program; by convention, a module variable whose value changes is named
between asterisks, as \code{*count*} is. Chapter~\ref{ch:syntax} covers
variables, and Chapter~\ref{ch:modules} covers modules, the namespaces
module variables live in (\dylanbook{chapters 4, 5 and 7}).

\section{Defining functions}\label{sec:tut-functions}

\index{define-method}
The usual way to define a function is \code{define-method}. It takes the
function's name, a list of parameters, and a body:

\begin{dylan}
(define-method square ((n <number>))
  (* n n))                      ; => square

(square 12)                     ; => 144
(square 1/2)                    ; => 1/4
\end{dylan}

The parameter \code{(n <number>)} says that \code{n} must be a number. The
class is called a \term{specializer}; if you leave it out and write just
\code{n}, the parameter accepts any object. A function returns the value of
the last form in its body; there is no \code{return} statement. Functions
can of course call other functions:

\begin{dylan}
(define-method hypotenuse ((a <real>) (b <real>))
  (sqrt (+ (square a) (square b))))

(hypotenuse 5 12)               ; => 13.0
\end{dylan}

\code{define-method} does a little more than define a function: it adds a
method to a \emph{generic function}, which can hold several methods for
different classes of argument. Section~\ref{sec:tut-generic} returns to
this, and it is the heart of Dylan. For now, think of
\code{define-method} as ``define a function''.

\paragraph{Anonymous functions.}\index{method}
The special form \code{method} makes a function without naming it. Such a
function is a value like any other: it can be called directly, stored in
a variable, or passed to another function.

\begin{dylan}
((method (x y) (* x y)) 6 7)    ; => 42

(define add1 (method (x) (+ x 1)))
(add1 41)                       ; => 42
\end{dylan}

\paragraph{Keyword and rest parameters.}\index{\#key}\index{\#rest}
After its required parameters, a parameter list may name \term{keyword
parameters} after \code{\#key}. The caller passes them by name, in any
order, and each may have a default value. A parameter after \code{\#rest}
collects any remaining arguments in a sequence.

\begin{dylan}
(define-method greet (name #key (greeting "Hello"))
  (concatenate greeting ", " name "!"))

(greet "Ada")                       ; => "Hello, Ada!"
(greet "Ada" greeting: "Goodbye")   ; => "Goodbye, Ada!"

(define-method sum-all (#rest numbers)
  (reduce + 0 numbers))

(sum-all 1 2 3 4 5)                 ; => 15
\end{dylan}

\paragraph{Returning several values.}\index{values}
\code{values} returns more than one value at once, and \code{bind}
receives them:

\begin{dylan}
(define-method smallest-and-largest ((s <sequence>))
  (values (reduce min (first s) s)
          (reduce max (first s) s)))

(smallest-and-largest '(3 1 4 1 5 9 2 6))   ; => 1, 9

(bind ((lo hi (smallest-and-largest #(10 -2 7))))
  (- hi lo))                                ; => 12
\end{dylan}

\paragraph{Closures.}\index{closure}
A function made inside another function can use the local variables
around it, and it keeps them for as long as it exists, even after the
function that made it has returned. Here \code{make-counter} returns a new
function each time it is called, and each of those functions has its own
private count:

\begin{dylan}
(define-method make-counter ()
  (bind ((n 0))
    (method ()
      (set! n (+ n 1))
      n)))

(define tick (make-counter))
(tick)                          ; => 1
(tick)                          ; => 2
(tick)                          ; => 3
\end{dylan}

Chapter~\ref{ch:functions} describes functions and parameter lists in
detail (\dylanbook{chapters 6 and 10}).

\section{Making decisions}\label{sec:tut-conditionals}

\index{if}
\code{if} takes a test, a form to evaluate when the test is true, and a
form to evaluate when it is false:

\begin{dylan}
(define-method sign-of ((n <real>))
  (if (negative? n)
      "negative"
      (if (zero? n) "zero" "positive")))

(sign-of -4)                    ; => "negative"
(sign-of 0)                     ; => "zero"
\end{dylan}

\begin{warnbox}
In Dylan, \code{\#f} is the \emph{only} false value. Zero, the empty
list and the empty string all count as true:
\begin{dylan}
(if 0 'true 'false)             ; => true
(if '() 'true 'false)           ; => true
\end{dylan}
\end{warnbox}

Each branch of an \code{if} is a single form. To do several things in a
branch, group them with \code{begin}, which evaluates its forms in order
and returns the value of the last.

\index{when}\index{unless}
\code{when} and \code{unless} are handy when there is only one branch.
\code{when} evaluates its body if the test is true, \code{unless} if it is
false; otherwise each returns \code{\#f}. Their bodies may contain several
forms.

\begin{dylan}
(when (> 3 2) "yes")            ; => "yes"
(when (> 2 3) "yes")            ; => #f
(unless (> 2 3) "no")           ; => "no"
\end{dylan}

\index{cond}
\code{cond} chooses among several cases. Each clause starts with a test;
the first clause whose test is true is evaluated, and the rest are skipped.
A final clause whose test is \code{else:} catches everything else.

\begin{dylan}
(define-method classify ((n <integer>))
  (cond ((< n 0)  "negative")
        ((= n 0)  "zero")
        ((< n 10) "small")
        (else:    "large")))

(classify 7)                    ; => "small"
(classify 42)                   ; => "large"
\end{dylan}

\index{case}
\code{case} compares one value against lists of constants, using
\code{id?}. The constants are not evaluated, so symbols in them need no
quote marks.

\begin{dylan}
(define-method day-kind (day)
  (case day
    ((saturday sunday) "weekend")
    ((monday tuesday wednesday thursday friday) "weekday")
    (else: "unknown")))

(day-kind 'sunday)              ; => "weekend"
(day-kind 'friday)              ; => "weekday"
(day-kind 'someday)             ; => "unknown"
\end{dylan}

Without an \code{else:} clause, a \code{case} with no matching clause
signals an error. Its relative \code{select} lets you choose the test that
compares the values; with \code{instance?} as the test it chooses by
class. Chapter~\ref{ch:control} describes it.

\index{and}\index{or}\index{not}
\code{and} and \code{or} evaluate their arguments from left to right and
stop as soon as the answer is known. \code{and} returns \code{\#f} or the
value of its last argument; \code{or} returns the first true value.
\code{not} turns \code{\#f} into \code{\#t} and every other value into
\code{\#f}.

\begin{dylan}
(and 1 2 3)                     ; => 3
(or #f 2 3)                     ; => 2
(not 0)                         ; => #f
\end{dylan}

Chapter~\ref{ch:control} covers all of these (\dylanbook{chapter 8}).

\section{Output and loops}\label{sec:tut-loops}

\paragraph{Printing.}\index{display}\index{print}\index{newline}
Before writing loops it helps to be able to print. Both implementations
have \code{display}, which writes a string or character as it is, and any
other object in its printed form; \code{newline}, which ends the line; and
\code{print}, which writes an object the way you would type it, strings in
quotes, followed by a newline.

\begin{dylan}
(begin
  (display "The answer is ")
  (display 42)
  (newline))
; prints: The answer is 42
\end{dylan}

\noindent
\code{format} does the same in one call, \code{(format \#t "The answer
is \textasciitilde a\textasciitilde\%" 42)}; Chapter~\ref{ch:printing}
describes it and the other printing functions.

\paragraph{for.}\index{for}
\code{for} is the general loop. Its first argument lists the loop
variables, each with an initial value and a form that computes its next
value. The second is a list of an end test and the forms whose value the
loop returns when the test becomes true. Any further forms are a body,
evaluated on each pass. This loop computes a factorial without any body
at all:

\begin{dylan}
(define-method factorial ((n <integer>))
  (for ((i n (- i 1))
        (product 1 (* product i)))
       ((<= i 0) product)))

(factorial 10)                  ; => 3628800
\end{dylan}

On each pass, all the variables are stepped together, using the values
from the previous pass.

\paragraph{dotimes.}\index{dotimes}
\code{dotimes} runs its body a fixed number of times, with a variable
counting from zero:

\begin{dylan}
(define-method sum-to ((n <integer>))
  (bind ((total 0))
    (dotimes (i (+ n 1))
      (set! total (+ total i)))
    total))

(sum-to 100)                    ; => 5050
\end{dylan}

\paragraph{for-each.}\index{for-each}
\code{for-each} steps through the elements of one or more collections.
Like \code{for}, it takes a list of variables, an end test with result
forms, and a body. The loop also ends, returning \code{\#f}, when any of
the collections runs out.

\begin{dylan}
(define-method first-even ((s <sequence>))
  (for-each ((x s))
            ((even? x) x)))

(first-even '(1 3 8 5 6))       ; => 8
(first-even '(1 3 5))           ; => #f
\end{dylan}

When there is no end test, write an empty list in its place. Here two
collections are stepped together, one of them an unbounded range of
integers:

\begin{dylan}
(for-each ((name '(ann bob cy))
           (i (range from: 1)))
          ()
  (display i) (display ". ") (display name) (newline))
; prints: 1. ann
;         2. bob
;         3. cy
\end{dylan}

\paragraph{while and until.}\index{while}\index{until}
\code{while} repeats its body as long as a test is true, and \code{until}
repeats it until a test becomes true. Both return \code{\#f}, so they are
used for their effects:

\begin{dylan}
(define-method collatz-steps ((n <integer>))
  (bind ((steps 0))
    (until (= n 1)
      (set! n (if (even? n) (/ n 2) (+ (* 3 n) 1)))
      (set! steps (+ steps 1)))
    steps))

(collatz-steps 27)              ; => 111
\end{dylan}

A method's parameters are ordinary local variables, so \code{set!} can
change them, as this one changes \code{n}.

In practice, many loops in Dylan are written not with these forms but
with the collection functions of Section~\ref{sec:tut-functional}.
Chapter~\ref{ch:control} covers the loops (\dylanbook{chapter 8}).

\section{Lists, vectors and strings}\label{sec:tut-sequences}

\index{sequences}
Lists, vectors and strings are all \term{sequences}: collections whose
elements are numbered from zero. They share most of their operations.

\paragraph{Lists.}\index{lists}
\code{list} makes a list of its arguments. \code{first}, \code{second}
and \code{third} take elements from the front; \code{element} takes any
element by its index; \code{size} counts them.

\begin{dylan}
(define primes (list 2 3 5 7 11))

(first primes)                  ; => 2
(second primes)                 ; => 3
(element primes 4)              ; => 11
(size primes)                   ; => 5
(last primes)                   ; => 11
\end{dylan}

A list is a chain of \term{pairs}. \code{cons} makes a new pair, putting
an element on the front of a list; \code{car} is a list's first element
and \code{cdr} the rest of it.

\begin{dylan}
(cons 1 (list 2 3))             ; => (1 2 3)
(car primes)                    ; => 2
(cdr primes)                    ; => (3 5 7 11)
(empty? '())                    ; => #t
\end{dylan}

Many operations make new sequences and leave their arguments alone:

\begin{dylan}
(reverse primes)                ; => (11 7 5 3 2)
(append primes '(13 17))        ; => (2 3 5 7 11 13 17)
(add primes 1)                  ; => (1 2 3 5 7 11)
primes                          ; => (2 3 5 7 11)
\end{dylan}

\code{add} adds an element wherever is most natural for the kind of
sequence, which for a list is the front.

\paragraph{Vectors.}\index{vectors}
A vector holds its elements in a block, so any element can be reached in
the same time. \code{vector} makes one from its arguments; \code{make}
makes one of a given size. \code{set!} with \code{element} changes an
element.

\begin{dylan}
(define scores (vector 10 20 30))

(element scores 1)              ; => 20
(set! (element scores 1) 99)    ; => 99
scores                          ; => #(10 99 30)

(make <vector> size: 3 fill: 0) ; => #(0 0 0)
\end{dylan}

The form \code{(set! (element scores 1) 99)} works because \code{set!} can
change more than variables: given a call such as \code{(element scores
1)}, it calls the matching \term{setter} function instead. You will see
the same thing with the slots of objects in Section~\ref{sec:tut-classes}.

A vector's size is fixed. A \code{<stretchy-vector>} grows as you add to
it with \code{add!}:

\begin{dylan}
(define names (make <stretchy-vector>))
(add! names 'ada)
(add! names 'grace)
(size names)                    ; => 2
(element names 1)               ; => grace
\end{dylan}

\paragraph{Strings.}\index{strings}
A string is a sequence of characters, so the sequence operations work on
strings too:

\begin{dylan}
(define lang "Dylan")

(size lang)                     ; => 5
(element lang 0)                ; => #\D
(concatenate lang " 1992")      ; => "Dylan 1992"
(as-uppercase lang)             ; => "DYLAN"
(reverse "abc")                 ; => "cba"
(copy-sequence "Hello, world" start: 7)   ; => "world"
\end{dylan}

\code{=} compares sequences element by element, while \code{id?} asks
whether two values are the same object. \code{as} converts between kinds
of sequence:

\begin{dylan}
(= "abc" "abc")                     ; => #t
(id? "abc" (copy-sequence "abc"))   ; => #f
(as <list> "abc")                   ; => (#\a #\b #\c)
(as <string> '(#\x #\y))            ; => "xy"
\end{dylan}

\paragraph{Ranges.}\index{ranges}
A \code{<range>} is a sequence of numbers in arithmetic progression. It
does not store its elements but computes them, so a range can be very
large or even unbounded, as the one in the \code{for-each} example of
Section~\ref{sec:tut-loops} was. \code{range} makes one: \code{from:} is
the first number, \code{up-to:} an upper bound that is not included,
\code{through:} a bound that is included, \code{by:} the step, and
\code{size:} the number of elements.

\begin{dylan}
(as <list> (range from: 1 up-to: 5))            ; => (1 2 3 4)
(as <list> (range from: 1 through: 5))          ; => (1 2 3 4 5)
(as <list> (range from: 10 through: 0 by: -5))  ; => (10 5 0)
(as <list> (range from: 0 size: 4))             ; => (0 1 2 3)
(reduce + 0 (range from: 1 through: 100))       ; => 5050
(member? 50 (range from: 0 by: 10))             ; => #t
\end{dylan}

\paragraph{Deques.}\index{deques}
A \code{<deque>}, a double-ended queue, grows and shrinks efficiently at
both ends. \code{push} and \code{pop} work at the front, and
\code{push-last} and \code{pop-last} at the back:

\begin{dylan}
(define line-up (make <deque>))
(push line-up 1)
(push-last line-up 2)
(push-last line-up 3)
(as <list> line-up)             ; => (1 2 3)
(pop line-up)                   ; => 1
(pop-last line-up)              ; => 3
(as <list> line-up)             ; => (2)
\end{dylan}

\paragraph{Sorting.}\index{sort}
\code{sort} returns a sorted copy of a sequence. Its \code{test:} keyword
argument is the function that compares two elements:

\begin{dylan}
(sort '(3 1 2) test: <)                     ; => (1 2 3)
(sort (vector "pear" "apple" "fig") test: <) ; => #("apple" "fig" "pear")
(sort '(3 1 2) test: >)                     ; => (3 2 1)
\end{dylan}

Chapter~\ref{ch:collections} describes every kind of sequence and its
operations (\dylanbook{chapter 12}).

\section{Tables}\label{sec:tut-tables}

\index{tables}
A \code{<table>} maps keys to values; other languages call it a hash
table, dictionary or map. Make one with \code{make}, store values with
\code{set!} and \code{element}, and read them back with \code{element}:

\begin{dylan}
(define ages (make <table>))

(set! (element ages 'alice) 34)
(set! (element ages 'bob) 27)
(set! (element ages 'carol) 41)

(element ages 'bob)                     ; => 27
(size ages)                             ; => 3
(element ages 'dave default: 'unknown)  ; => unknown
\end{dylan}

Asking for a key that is not in the table is an error unless you supply
\code{default:}, whose value is returned instead. \code{remove-key!}
removes a key. A table is a collection whose elements are its values, so
the collection operations apply to it; \code{key-sequence} gives its keys.

\begin{dylan}
(reduce + 0 ages)                       ; => 102
(remove-key! ages 'bob)
(size ages)                             ; => 2
\end{dylan}

Here is a common use: counting how often each word occurs.

\begin{dylan}
(define word-counts (make <table>))

(for-each ((w '(the cat saw the dog and the bird)))
          ()
  (set! (element word-counts w)
        (+ 1 (element word-counts w default: 0))))

(element word-counts 'the)              ; => 3
(element word-counts 'cat)              ; => 1
\end{dylan}

The order of a table's elements and keys is not defined, so sort them when
the order matters. Tables are described with the other collections in
Chapter~\ref{ch:collections}.

\section{Classes and objects}\label{sec:tut-classes}

\index{classes}\index{define-class}
\code{define-class} defines a class. It takes the class's name, a list of
superclasses, and a description of each \term{slot}. Every class
ultimately inherits from \code{<object>}.

\begin{dylan}
(define-class <point> (<object>)
  (x init-value: 0 init-keyword: x:)
  (y init-value: 0 init-keyword: y:))
\end{dylan}

Each slot here has a name, a default value given by \code{init-value:},
and a keyword given by \code{init-keyword:} that sets the slot when an
object is made. \code{make} makes an \term{instance} of a class:

\index{make}
\begin{dylan}
(define pt (make <point> x: 3 y: 4))

(x pt)                          ; => 3
(y pt)                          ; => 4
(y (make <point>))              ; => 0
\end{dylan}

Defining the class also defined a \term{getter} function for each slot,
named after the slot, which is why \code{(x pt)} reads the slot \code{x}.
It defined a setter too, which \code{set!} uses:

\begin{dylan}
(set! (x pt) 10)                ; => 10
(x pt)                          ; => 10
\end{dylan}

\code{instance?} asks whether an object belongs to a class, and
\code{object-class} answers the object's class:

\begin{dylan}
(instance? pt <point>)          ; => #t
(instance? 5 <point>)           ; => #f
(object-class pt)
\end{dylan}

\code{(object-class pt)} prints as \code{\{the class <point>\}}, and
\code{pt} itself, like any object without a readable printed form, as
\code{\{instance of <point>\}}. Chapter~\ref{ch:printing} explains how
to give your own classes a printed form.

\paragraph{Inheritance.} A class inherits the slots of its superclasses
and adds its own. Here is a small family of shapes, which the following
sections use:

\begin{dylan}
(define-class <shape> (<object>)
  (name init-value: "shape" init-keyword: name:))

(define-class <circle> (<shape>)
  (radius required-init-keyword: radius:))

(define-class <rect> (<shape>)
  (width  init-keyword: width:)
  (height init-keyword: height:))
\end{dylan}

\code{required-init-keyword:} means that \code{make} must be given that
keyword; making a \code{<circle>} without a \code{radius:} is an error.
Instances of the subclasses have the inherited \code{name} slot as well
as their own, and they are instances of \code{<shape>} as well as of their
own classes:

\begin{dylan}
(define wheel (make <circle> radius: 2 name: "wheel"))

(name wheel)                    ; => "wheel"
(radius wheel)                  ; => 2
(instance? wheel <shape>)       ; => #t
(subclass? <circle> <shape>)    ; => #t
\end{dylan}

\paragraph{Multiple inheritance.} A class may have more than one
superclass, and it then inherits from all of them:

\begin{dylan}
(define-class <colored> (<object>)
  (color init-value: 'black init-keyword: color:))

(define-class <colored-circle> (<circle> <colored>))

(define dot (make <colored-circle> radius: 1 color: 'red))
(list (color dot) (radius dot) (name dot))   ; => (red 1 "shape")
\end{dylan}

\paragraph{Initializing objects.}\index{initialize}
A slot's initial value need not be a constant. \code{init-function:}
gives a function of no arguments that is called to compute the value each
time an object is made. This matters when each object needs its own
mutable value, such as a fresh collection, or a value computed afresh,
such as a serial number:

\begin{dylan}
(define *next-id* 0)

(define-class <ticket> (<object>)
  (id    init-function: (method ()
                          (set! *next-id* (+ *next-id* 1))
                          *next-id*))
  (title init-value: "untitled" init-keyword: title:)
  (notes init-function: (method () (make <stretchy-vector>))))

(id (make <ticket>))                         ; => 1
(id (make <ticket> title: "printer jam"))    ; => 2
\end{dylan}

Had \code{notes} used \code{init-value:} instead, every ticket would share
one and the same stretchy vector.

For initialization that slot options cannot express, add a method to the
generic function \code{initialize}, which \code{make} calls on each new
object with the keyword arguments it was given. Call \code{next-method}
first, so that the slots are filled in the usual way, and then do the
rest. Here a rectangle made without a height becomes a square;
\code{slot-initialized?} asks whether a slot has a value yet:

\begin{dylan}
(define-class <rectangle> (<object>)
  (width  init-keyword: width:)
  (height init-keyword: height:))

(define-method initialize ((r <rectangle>) #key)
  (next-method)
  (unless (slot-initialized? r height)
    (set! (height r) (width r))))

(height (make <rectangle> width: 4))             ; => 4
(height (make <rectangle> width: 4 height: 3))   ; => 3
\end{dylan}

\begin{warnbox}
A slot's getter is an ordinary function in the module, so a local
variable or parameter with the same name hides it. Inside a method with a
parameter called \code{name}, the call \code{(name s)} would try to call
the parameter rather than the getter. Give parameters names that differ
from your slots.
\end{warnbox}

Chapter~\ref{ch:classes} describes classes, slot options and the
initialization protocol (\dylanbook{chapters 6 and 11}).

\section{Generic functions and multiple dispatch}\label{sec:tut-generic}

\index{generic function}\index{multiple dispatch}
Every \code{define-method} you have written so far has added a
\term{method} to a \term{generic function}. A generic function can hold
many methods, each for a different class of argument. When the generic
function is called, it chooses the most specific method whose
specializers match the arguments and runs it.

\code{define-generic-function} declares a generic function and its
parameters. It is optional---the first \code{define-method} creates the
generic function if there is none---but it documents the protocol in one
place:

\begin{dylan}
(define-generic-function area (shape))

(define-method area ((c <circle>))
  (* 3.14159 (radius c) (radius c)))

(define-method area ((r <rect>))
  (* (width r) (height r)))

(area (make <rect> width: 3 height: 5))     ; => 15
(area wheel)                                ; => 12.56636
\end{dylan}

Because \code{area} works on any shape, functions that use it work on any
shape too:

\begin{dylan}
(define some-shapes
  (list (make <circle> radius: 1) (make <rect> width: 2 height: 2)))

(map area some-shapes)          ; => (3.14159 4)
\end{dylan}

\paragraph{The most specific method wins.} Methods may overlap. When more
than one applies, the one whose specializers are more specific is chosen.
A method on \code{<shape>} serves every shape that has no method of its
own:

\begin{dylan}
(define-method describe ((s <shape>))
  (concatenate "a shape called " (name s)))

(describe (make <rect> width: 1 height: 1))   ; => "a shape called shape"
\end{dylan}

\paragraph{Dispatch on several arguments.} A method can be specialized on
any or all of its required arguments, and the choice takes all of them
into account. This is \term{multiple dispatch}, and it is what makes
generic functions more than methods attached to a single class:

\begin{dylan}
(define-method collide ((a <circle>) (b <circle>))
  "two circles")

(define-method collide ((a <circle>) (b <rect>))
  "a circle and a rectangle")

(define-method collide ((a <shape>) (b <shape>))
  "two shapes")

(collide wheel (make <rect> width: 1 height: 1))  ; => "a circle and a rectangle"
(collide (make <rect> width: 1 height: 1) wheel)  ; => "two shapes"
(collide wheel wheel)                             ; => "two circles"
\end{dylan}

\paragraph{Specializing on a single object.}\index{singleton}
A specializer can also be one particular object, written with
\code{singleton}. Such a method applies only when the argument is that
object, and it is more specific than any method on a class. This gives a
neat way to write the base case of a recursion:

\begin{dylan}
(define-method fact ((n <integer>))
  (* n (fact (- n 1))))

(define-method fact ((n (singleton 0)))
  1)

(fact 5)                        ; => 120
\end{dylan}

Singletons also suit symbols that act as commands or names:

\begin{dylan}
(define-method speak ((animal (singleton 'dog))) "woof")
(define-method speak ((animal (singleton 'cat))) "meow")
(define-method speak ((animal <symbol>)) "...")

(speak 'dog)                    ; => "woof"
(speak 'fish)                   ; => "..."
\end{dylan}

Built-in classes such as \code{<integer>}, \code{<string>} and
\code{<symbol>} are classes like any other, so you can add methods for them
to your own generic functions. Chapter~\ref{ch:functions} describes
generic functions and how methods are ordered (\dylanbook{chapters 6 and
10}).

\section{Calling the next method}\label{sec:tut-next-method}

\index{next-method}
Often a method for a subclass wants to do what the method for the
superclass does, and a little more. Inside a method, \code{next-method}
calls the next most specific method, the one that would have run if this
one did not exist. Called with no arguments, it passes along the original
arguments:

\begin{dylan}
(define-method describe ((c <circle>))
  (concatenate "a circle, " (next-method)))

(describe wheel)                ; => "a circle, a shape called wheel"
\end{dylan}

Here the method on \code{<circle>} runs first and calls the method on
\code{<shape>} for the rest of the description. \code{next-method} may
also be given arguments, which replace the original ones; the bank
example in Section~\ref{sec:tut-program} uses this to change an amount
before passing it on. Chapter~\ref{ch:functions} has the details.

\section{Collections and functional operations}\label{sec:tut-functional}

\index{map}\index{reduce}\index{choose}
Because functions are values, much of Dylan's power comes from functions
that take other functions as arguments. These work on every kind of
collection.

\code{map} applies a function to each element and collects the results in
a new collection of the same kind as the first argument. Given several
collections, it takes one element from each. \code{map-as} names the kind
of collection to return.

\begin{dylan}
(map (method (x) (* x x)) '(1 2 3 4))   ; => (1 4 9 16)
(map + '(1 2 3) '(10 20 30))            ; => (11 22 33)
(map-as <vector> square '(1 2 3))       ; => #(1 4 9)
\end{dylan}

\code{reduce} combines the elements with a function of two arguments,
starting from an initial value; \code{reduce1} uses the first element as
the initial value.

\begin{dylan}
(reduce + 0 '(1 2 3 4 5))       ; => 15
(reduce1 max '(3 9 2))          ; => 9
\end{dylan}

\code{choose} keeps the elements that satisfy a predicate.
\code{any?} and \code{every?} ask whether some or all of them do.
\code{find-key} answers the key, or index, of the first element that
satisfies a predicate.

\begin{dylan}
(choose even? '(1 2 3 4 5 6))   ; => (2 4 6)
(any? negative? '(3 -1 4))      ; => #t
(every? positive? '(3 -1 4))    ; => #f
(find-key #(5 8 13) even?)      ; => 1
\end{dylan}

\code{do} calls a function on each element just for its effect, and
returns \code{\#f}:

\begin{dylan}
(do (method (s) (display (area s)) (newline)) some-shapes)
; prints: 3.14159
;         4
\end{dylan}

\paragraph{Making functions from functions.}\index{curry}\index{compose}
\index{complement}
\code{curry} fixes the first arguments of a function. \code{compose}
chains functions, the rightmost applied first. \code{complement} turns a
predicate into its opposite. \code{apply} calls a function with arguments
taken from a sequence.

\begin{dylan}
(define add10 (curry + 10))
(add10 5)                                ; => 15
(map (curry * 2) '(1 2 3))               ; => (2 4 6)

(define sqrt-of-abs (compose sqrt abs))
(sqrt-of-abs -16)                        ; => 4.0

(choose (complement even?) '(1 2 3 4 5)) ; => (1 3 5)

(apply max '(4 7 1))                     ; => 7
\end{dylan}

Together these replace many explicit loops. The total area of a list of
shapes, for instance, is one line:

\begin{dylan}
(reduce + 0 (map area some-shapes))      ; => 7.14159
\end{dylan}

Chapter~\ref{ch:collections} lists the collection operations and
Chapter~\ref{ch:functional} the functional operations
(\dylanbook{chapters 12 and 17}).

\section{Leaving early with bind-exit}\label{sec:tut-exits}

\index{bind-exit}
Sometimes you want to stop a computation partway through and return a
value at once, for example when a search finds what it was looking for.
\code{bind-exit} binds a name to an \term{exit procedure}. Calling that
procedure, from anywhere inside the \code{bind-exit}, however deeply
nested, ends the \code{bind-exit} immediately and makes its argument the
value of the whole form.

\begin{dylan}
(define-method find-first ((pred <function>) (s <sequence>))
  (bind-exit (return)
    (do (method (x)
          (when (pred x) (return x)))
        s)
    #f))

(find-first even? '(1 3 4 5 6))     ; => 4
(find-first even? '(1 3 5))         ; => #f
\end{dylan}

The exit procedure is called from inside the function passed to
\code{do}, yet it leaves \code{find-first} entirely. If it is never called,
\code{bind-exit} returns the value of its last form, here \code{\#f}.

\index{unwind-protect}
\code{unwind-protect} guarantees that some cleanup forms run, however its
first form is left, whether normally or by an exit:

\begin{dylan}
(bind-exit (leave)
  (unwind-protect
    (leave 'escaped)
    (display "cleanup runs anyway")
    (newline)))
; prints: cleanup runs anyway
; => escaped
\end{dylan}

Chapter~\ref{ch:control} describes non-local exits
(\dylanbook{chapter 8}).

\section{Conditions and errors}\label{sec:tut-conditions}

\index{conditions}\index{error}
When something goes wrong, a Dylan program \term{signals} a
\term{condition}, an object that describes what happened. \code{error}
signals an error. In its simplest form it takes a message, which may
contain \code{\textasciitilde{}a} where each further argument is to be
inserted:

\begin{dylan}
(define-method safe-divide ((a <number>) (b <number>))
  (if (zero? b)
      (error "cannot divide ~a by zero" a)
      (/ a b)))

(safe-divide 10 4)              ; => 5/2
\end{dylan}

If nothing handles the error, the interactive implementation enters the
debugger, as Section~\ref{sec:start-rep} showed, and a compiled program
prints the message and exits.

\paragraph{Handling errors.}\index{handler-case}
\code{handler-case} evaluates a form and, if a condition of a given class
is signaled while it runs, abandons the form and evaluates the matching
clause instead. Each clause begins with a list of the class and, after
\code{condition:}, an optional name for the condition object:

\begin{dylan}
(handler-case (safe-divide 1 0)
  ((<error> condition: c)
   (list 'caught (condition-format-string c))))
; => (caught "cannot divide ~a by zero")
\end{dylan}

An error signaled by \code{error} with a message is an instance of
\code{<simple-error>}. The function \code{condition-format-string}
returns its message, and \code{condition-format-arguments} returns the
other arguments in a list. Errors that the system itself signals are
conditions too, and \code{handler-case} catches them the same way:

\begin{dylan}
(handler-case (element (vector 1 2) 5)
  ((<error>) "no such element"))    ; => "no such element"

(handler-case (+ 1 'a)
  ((<error>) "not a number"))       ; => "not a number"
\end{dylan}

\paragraph{Your own conditions.} Conditions are objects, and their classes
are ordinary classes. To signal something more specific than a message,
define a subclass of \code{<error>} with slots that describe the problem,
and pass an instance of it to \code{error}:

\begin{dylan}
(define-class <out-of-stock> (<error>)
  (item  required-init-keyword: item:)
  (short required-init-keyword: short:))

(define-method take-from-stock ((on-hand <integer>) (wanted <integer>))
  (if (> wanted on-hand)
      (error (make <out-of-stock> item: 'widget short: (- wanted on-hand)))
      (- on-hand wanted)))

(take-from-stock 10 4)          ; => 6

(handler-case (take-from-stock 10 15)
  ((<out-of-stock> condition: e)
   (list 'short-by (short e))))  ; => (short-by 5)
\end{dylan}

A \code{handler-case} may have several clauses; they are tried in order,
so put the more specific classes first.

\paragraph{Signals that are not errors.}\index{signal}\index{warnings}
\code{signal} signals a condition that a program may ignore. Given a
message, it signals a \code{<simple-warning>}. If no handler takes it,
\code{signal} prints the warning and returns \code{\#f}, and the program
carries on:

\begin{dylan}
(handler-case (signal "low on ~a" 'paper)
  ((<warning> condition: w)
   (condition-format-arguments w)))   ; => (paper)

(signal "nobody handles this")    ; => #f
; prints: WARNING: nobody handles this          (interactive)
;         Dylan warning: nobody handles this    (compiled)
\end{dylan}

Dylan's condition system can do much more: \code{handler-bind}
establishes a handler that runs \emph{without} first leaving the
signaling code, so it can return a value to it, decline, or choose among
\term{restarts}. Chapter~\ref{ch:conditions} explains it all
(\dylanbook{chapter 13}).

\section{A complete program}\label{sec:tut-program}

This section builds a small but complete program, a model of a bank, one
piece at a time. It brings together classes, generic functions,
\code{next-method}, tables, the collection functions and conditions. Type
the pieces at the loop as you go, or put them all in a file called
\file{bank.dyl}; the next section compiles that file.

\paragraph{Accounts.} Every account has an owner and a balance. A savings
account also earns interest at some rate, and a checking account charges
a fee on each withdrawal. Balances are whole numbers of dollars, and the
interest rate is an exact ratio, so no money is lost to rounding until we
choose to round.

\begin{dylan}
(define-class <account> (<object>)
  (owner   required-init-keyword: owner:)
  (balance init-value: 0 init-keyword: balance:))

(define-class <savings-account> (<account>)
  (rate init-value: 1/50 init-keyword: rate:))

(define-class <checking-account> (<account>)
  (fee init-value: 1 init-keyword: fee:))
\end{dylan}

\paragraph{An error class.} Withdrawing more than the balance is an
error, and the condition records which account and how much was asked
for. The slot names are chosen so that they do not clash with the names of
parameters used elsewhere (see the warning in
Section~\ref{sec:tut-classes}).

\begin{dylan}
(define-class <insufficient-funds> (<error>)
  (refused-account required-init-keyword: account:)
  (refused-amount  required-init-keyword: amount:))
\end{dylan}

\paragraph{The operations.} Three generic functions make up the protocol
of an account. The basic methods are on \code{<account>}:

\begin{dylan}
(define-generic-function deposit! (account amount))
(define-generic-function withdraw! (account amount))
(define-generic-function month-end! (account))

(define-method deposit! ((a <account>) (amount <integer>))
  (set! (balance a) (+ (balance a) amount)))

(define-method withdraw! ((a <account>) (amount <integer>))
  (when (> amount (balance a))
    (error (make <insufficient-funds> account: a amount: amount)))
  (set! (balance a) (- (balance a) amount)))

(define-method month-end! ((a <account>))
  (balance a))
\end{dylan}

Try them on a plain account:

\begin{dylan}
(define acct (make <account> owner: 'dave balance: 100))

(deposit! acct 25)              ; => 125
(withdraw! acct 40)             ; => 85
(balance acct)                  ; => 85
\end{dylan}

\paragraph{Specializing with next-method.} The other kinds of account
differ in small ways, which is exactly what \code{next-method} is for. A
checking account's withdrawal is an ordinary withdrawal of the amount
\emph{plus the fee}, so its method calls the next method with a changed
argument. A savings account adds its interest at the end of the month and
then does whatever any account does.

\begin{dylan}
(define-method withdraw! ((a <checking-account>) (amount <integer>))
  (next-method a (+ amount (fee a))))

(define-method month-end! ((a <savings-account>))
  (deposit! a (round (* (balance a) (rate a))))
  (next-method))
\end{dylan}

\code{round} returns two values, the nearest integer and the remainder,
but \code{deposit!} uses only the first.

\begin{dylan}
(define chk (make <checking-account> owner: 'erin balance: 100 fee: 2))
(withdraw! chk 10)              ; => 88

(define sav (make <savings-account> owner: 'fred balance: 500))
(month-end! sav)                ; => 510
\end{dylan}

Because the withdrawal for a checking account goes through the method on
\code{<account>}, it also gets that method's check for insufficient
funds, fee included:

\begin{dylan}
(handler-case (withdraw! chk 87)
  ((<insufficient-funds> condition: c)
   (refused-amount c)))         ; => 89
\end{dylan}

\paragraph{The bank.} The bank keeps its accounts in a table, keyed by
the owner's name. \code{open-account} takes the class of account to open,
so one function opens every kind. It passes any further keyword arguments
on to \code{make} with \code{apply}.

\begin{dylan}
(define *accounts* (make <table>))

(define-method open-account ((class <class>) (who <symbol>) #rest options)
  (bind ((a (apply make class owner: who options)))
    (set! (element *accounts* who) a)
    a))

(define-method lookup ((who <symbol>))
  (element *accounts* who default: #f))
\end{dylan}

\paragraph{Handling a refusal.} \code{try-withdraw} attempts a withdrawal
and reports a refusal instead of letting the error escape. It answers
whether the withdrawal was made.

\begin{dylan}
(define-method try-withdraw ((who <symbol>) (amount <integer>))
  (handler-case (begin (withdraw! (lookup who) amount) #t)
    ((<insufficient-funds> condition: c)
     (display "refused: ")
     (display (owner (refused-account c)))
     (display " cannot withdraw ")
     (display (refused-amount c))
     (newline)
     #f)))
\end{dylan}

\paragraph{Reports.} The total of all deposits maps \code{balance} over
the table and adds up the results. The report lists the accounts in
alphabetical order of owner; since the table's keys come in no particular
order, it sorts them, comparing the symbols as strings.

\begin{dylan}
(define-method total-deposits ()
  (reduce + 0 (map balance *accounts*)))

(define-method report ()
  (for-each ((who (sort (key-sequence *accounts*)
                        test: (method (a b)
                                (< (as <string> a) (as <string> b))))))
            ()
    (display who)
    (display ": ")
    (display (balance (lookup who)))
    (newline)))
\end{dylan}

\paragraph{A month at the bank.} Finally, the program itself: open three
accounts, make some transactions, close the month, and report.
\code{(do month-end! *accounts*)} runs the month's end for every account,
each with its own method.

\begin{dylan}
(open-account <savings-account> 'alice balance: 1000)
(open-account <checking-account> 'bob balance: 200 fee: 2)
(open-account <account> 'carol)

(deposit! (lookup 'carol) 50)
(withdraw! (lookup 'bob) 48)
(try-withdraw 'carol 80)
(do month-end! *accounts*)
(report)
(display "total: ")
(display (total-deposits))
(newline)
\end{dylan}

At the loop, each form's value appears after it as usual, mixed with the
output. The output itself is this:

\begin{shell}
refused: carol cannot withdraw 80
alice: 1020
bob: 150
carol: 50
total: 1220
\end{shell}

Alice earned two percent interest on 1000; Bob's withdrawal of 48 cost
him 50 with the fee; and Carol's attempt to withdraw 80 from her 50 was
refused and reported, without stopping the program.

\section{Compiling the program}\label{sec:tut-compile}

\index{dylanc}
The bank is a complete program, and \cmd{dylanc} can turn it into an
executable. Put the definitions and the final forms of the previous
section---everything except the small experiments with \code{acct},
\code{chk} and \code{sav}---into \file{bank.dyl}, in the order they were
given, and compile it:

\begin{shell}
$ bin/dylanc bank.dyl
$ ./bank
refused: carol cannot withdraw 80
alice: 1020
bob: 150
carol: 50
total: 1220
\end{shell}

The output is the same as at the loop, because the program uses only
\code{display} and \code{newline} for its output, which both
implementations have. Nothing in the program changes a class or adds a
method while it runs, so it is a static program, and the compiler accepts
it unchanged.

You can also run the same file in the interactive implementation without
typing it, either by feeding it to the loop or by compiling it to Lisp, as
Chapter~\ref{ch:start} showed:

\begin{shell}
$ (cat bank.dyl; echo dylan:done) | \
    sbcl --non-interactive --eval '(load "src/load-dylan.lisp")' \
         --eval '(dylan:rep)'
\end{shell}

This is the usual way of working with the two implementations: develop
and test a piece at a time at the loop, where every definition can be
changed and tried at once, then compile the finished program into a fast
executable that needs neither SBCL nor Dylan to run.

\begin{tipbox}
The compiler warns about any name a program uses but never defines.
Building with \cmd{dylanc} from time to time, even while you are still
working at the loop, is a quick way to catch misspelled names.
\end{tipbox}

\section{Where to go from here}\label{sec:tut-next}

You have now seen most of the language at least once. Part~II covers each
topic thoroughly, in the order of this tutorial: syntax and variables
(Chapter~\ref{ch:syntax}), control flow (Chapter~\ref{ch:control}),
functions and generic functions (Chapter~\ref{ch:functions}), classes
(Chapter~\ref{ch:classes}), numbers (Chapter~\ref{ch:numbers}),
collections (Chapter~\ref{ch:collections}), functional operations
(Chapter~\ref{ch:functional}), conditions (Chapter~\ref{ch:conditions}),
modules (Chapter~\ref{ch:modules}), and printing
(Chapter~\ref{ch:printing}).

Part~III describes the two implementations as tools. Read
Chapter~\ref{ch:dylanc} before building larger programs with the compiler,
and Chapter~\ref{ch:extensions} for its macros and threads.
Appendix~\ref{app:differences} lists everything in which this system
differs from the 1992 manual.

The 1992 manual itself is worth reading alongside this book. It is short,
clearly written, and full of examples, and its annotated copy in
\file{manual/Dylan-1992/} links each chapter to the design notes that
later changed it.
